\documentclass[11pt]{article}
\usepackage[a4paper,margin=2.0cm]{geometry}
\usepackage{amsmath,amssymb,booktabs,graphicx}
\usepackage[T1]{fontenc}
\usepackage{hyperref}
% generated by qbscreen/make_numbers_v6.py — do not edit
\newcommand{\AmineErr}{0.05}
\newcommand{\AnisoCellMax}{7.4}
\newcommand{\AnisoRowMax}{3.11}
\newcommand{\AnisoRowMin}{0.76}
\newcommand{\BornRadius}{2.81}
\newcommand{\CoulHi}{0.05}
\newcommand{\CoulLo}{0.03}
\newcommand{\Dcontact}{-843}
\newcommand{\DeltaMAOhi}{1.44}
\newcommand{\DeltaMAOlo}{1.41}
\newcommand{\DeltaMAOsub}{0.91}
\newcommand{\DeltaMAOsubHi}{0.99}
\newcommand{\DeltaMAOsubLo}{0.86}
\newcommand{\DeltaMaxHi}{0.82}
\newcommand{\DeltaMaxLo}{0.58}
\newcommand{\DeltaMaxPHi}{0.86}
\newcommand{\DeltaMaxPLo}{0.62}
\newcommand{\DeltaSiteHi}{2.00}
\newcommand{\DeltaSiteLo}{0.49}
\newcommand{\DeltaToThr}{0.41}
\newcommand{\EamExp}{1.30}
\newcommand{\Econtact}{-26}
\newcommand{\EflMAOA}{-159}
\newcommand{\EflMAOB}{-167}
\newcommand{\EmFreeA}{-211}
\newcommand{\EoneCap}{0.29}
\newcommand{\EoneShiftCap}{0.45}
\newcommand{\EsqredA}{-262}
\newcommand{\EsqredB}{-275}
\newcommand{\FkPhundred}{0.023}
\newcommand{\FkPlow}{14}
\newcommand{\FkPthousand}{2.6\times10^{-4}}
\newcommand{\FkPtop}{1.2\times10^{-7}}
\newcommand{\GdefOne}{0.25}
\newcommand{\GdefOneHi}{0.61}
\newcommand{\GsiteAniso}{0.21}
\newcommand{\GsiteRef}{0.06}
\newcommand{\MedOneHi}{0.04}
\newcommand{\MedOneLo}{-0.37}
\newcommand{\MedTwoHi}{0.56}
\newcommand{\MedTwoLo}{0.01}
\newcommand{\SubShift}{0.5}
\newcommand{\ThrHundredthOne}{0.53}
\newcommand{\ThrHundredthOneRelax}{0.53}
\newcommand{\ThrTenthOne}{0.47}
\newcommand{\ThrTenthOneAniso}{0.50}
\newcommand{\ThrTenthOneRelax}{0.47}
\newcommand{\ThrTenthOneSens}{0.50}
\newcommand{\ThrTenthThousand}{0.28}
\newcommand{\ThrTenthThousandAniso}{0.31}
\newcommand{\ThrTenthThousandRelax}{0.28}
\newcommand{\TtwoCellMax}{4.94}
\newcommand{\TtwoCellMin}{0.11}
\newcommand{\TtwoRowMax}{0.99}
\newcommand{\TtwoRowMin}{0.84}
\newcommand{\XTtwoMax}{4.78}
\newcommand{\XTtwoMin}{0.42}
\newcommand{\XanisoHMax}{2.2}
\newcommand{\XanisoMax}{3.1}
\newcommand{\XanisoNMax}{1.6}
\newcommand{\Xsens}{3.1}
\newcommand{\aBenzOne}{257}
\newcommand{\aBenzTwo}{196}
\newcommand{\aTrpOne}{60}
\newcommand{\aTrpTwo}{56}
\newcommand{\convMaxRatio}{1.018}
\newcommand{\cryDcoup}{-11.4}
\newcommand{\cryMax}{0.63}
\newcommand{\gridStep}{0.03}
\newcommand{\kTln}{0.037}
\newcommand{\maxBoundRatio}{0.9986}
\newcommand{\nAnisoCells}{30}
\newcommand{\nAnisoEdge}{0}
\newcommand{\nBoundTest}{300}
\newcommand{\nCells}{195}
\newcommand{\nCry}{21}
\newcommand{\nEdge}{0}
\newcommand{\nScan}{112}
\newcommand{\nScanValid}{69}
\newcommand{\nSuppRaised}{151}
\newcommand{\nTtwoCells}{90}
\newcommand{\offDelta}{0.46}
\newcommand{\offG}{0.0965}
\newcommand{\offKP}{31.6}
\newcommand{\offKShi}{379}
\newcommand{\offKSlo}{126}
\newcommand{\offMax}{0.0924}
\newcommand{\polishGain}{156}
\newcommand{\rowMaxChange}{0.0}
\newcommand{\sMax}{3}
\newcommand{\sharpRatio}{0.999996}
\newcommand{\suppMaxRatio}{3.6}
\newcommand{\tContactMax}{106}
\newcommand{\tContactMed}{53}

\renewcommand{\thetable}{S\arabic{table}}
\renewcommand{\thesection}{S\arabic{section}}
\renewcommand{\theequation}{S\arabic{equation}}

\title{Supporting Information:\\ Catalytic competence bounds the decay rates of
a thermally formed radical pair}
\author{Hikaru Wakaura and Taiki Tanimae}
\date{}

\begin{document}
\maketitle

\section{Software and environment}
Spin dynamics: \texttt{qbscreen} (this work) with NumPy and SciPy; coherent
yields from the Sylvester equation
$H_{\rm eff}X - XH_{\rm eff}^\dagger = -i\rho_0$ with
$H_{\rm eff} = 2\pi H - \tfrac{i}{2}(k_SP_S + k_P)$ and $X = \int_0^\infty\rho\,dt$,
solved by the Bartels--Stewart method;\cite{Bartels1972}
yields with electron spin relaxation from the Liouville-space generator. The two
agree to $10^{-8}$ without relaxation, and every yield pair sums to one.
Electronic structure: PySCF~2.12.1\cite{Sun2020} (B3LYP\cite{Becke1993,Lee1988};
def2-SVP for couplings and energies, def2-TZVP\cite{Weigend2005} for hyperfine
couplings; RI-J; PCM), geometries from GFN2-xTB~6.7.1.\cite{Bannwarth2019}
Every table below is generated from the stored results by
\url{qbscreen/make_si_v6.py}.

\section{The inequality $\Phi_P\le4k_P/k_S$ and the dead-end bound}
Let $\mathcal L$ be the generator of eq~2 of the main text, with $\mathcal D$
trace preserving, generating a positivity-preserving evolution (for example a
trace-preserving GKSL generator\cite{Gorini1976,Lindblad1976}) and unital,
$\mathcal D(I)=0$. Then $-\mathcal L(I) = k_SP_S + k_PI$. The rest of the
generator, $\mathcal L-\mathcal D$, generates
$\rho\mapsto e^{-iH_{\rm eff}t}\rho\,e^{iH_{\rm eff}^\dagger t}$ with
$H_{\rm eff} = 2\pi H - \tfrac i2(k_SP_S+k_PI)$, which is positive; by the
Lie--Trotter product formula,
$e^{\mathcal Lt} = \lim_{n\to\infty}(e^{\mathcal Dt/n}e^{(\mathcal L-\mathcal D)t/n})^n$
is a limit of positive maps and hence positive. For $k_P>0$ the operator
$\mathcal R = (-\mathcal L)^{-1} = \int_0^\infty e^{\mathcal Lt}\,dt$ therefore
exists and is positive, because the trace of $e^{\mathcal Lt}(\rho)$ decays at
least as $e^{-k_Pt}$ for positive $\rho$.
Applying $\mathcal R$ to the identity above gives
\begin{equation}
k_S\,\mathcal R(P_S) = I - k_P\,\mathcal R(I) \preceq I ,
\end{equation}
so ${\rm Tr}\,\mathcal R(P_S)\le 4N/k_S$ with $4N$ the dimension of the spin
space, and $\Phi_P = (k_P/N)\,{\rm Tr}\,\mathcal R(P_S)\le4k_P/k_S$.

The factor of four cannot be reduced for the class. With $H=0$, no nuclei,
$k_S = 1$ and $k_P = 10^{-6}~\mu$s$^{-1}$, the fully depolarizing generator
$\mathcal D(\rho) = \gamma({\rm Tr}\rho\,I/4-\rho)$, which is trace preserving
and unital, gives $\Phi_Pk_S/4k_P = 0.25$ at $\gamma=0$ and \sharpRatio\ for
$\gamma\ge10^2~\mu$s$^{-1}$ (\url{competence_sensitivity.sharpness_example};
\url{calibration_results/phi_bound_sharpness.json}). Trace preservation is not
needed for $\Phi_P\le4k_P/k_S$ but is needed for $\Phi_P\le1$: the unital,
completely positive but population-creating generator
$\mathcal D(\rho) = \gamma(V\rho V^\dagger-\tfrac12\{VV^\dagger,\rho\})$ with
$V = |T_0\rangle\langle S|$ gives $\Phi_P = 150/101$ at $k_S = 0.01$,
$k_P = \gamma = 1~\mu$s$^{-1}$ (\url{competence_sensitivity.trace_counterexample}).

For a dead-end pair ($k_P=0$) the resolvent need not exist: without nuclei and
with $H=0$, for example, the $16\times16$ generator has rank 7, because triplet
states that never reach the singlet are stationary. Integrating
$\frac{d}{du}e^{\mathcal L_0u}(I) = -k_S\,e^{\mathcal L_0u}(P_S)$ instead gives
\begin{equation}
k_S\int_0^t e^{\mathcal L_0u}(P_S)\,du = I - e^{\mathcal L_0t}(I) \preceq I
\end{equation}
for all $t$, since $e^{\mathcal L_0t}(I)\succeq0$. The population built up from a
constant singlet source $k_f[{\rm E\cdot S}]P_S/N$ is therefore at most
$4k_f[{\rm E\cdot S}]/k_S = e^{-\Delta/k_BT}[{\rm E\cdot S}]$ at all times.

The function \url{random_bound_test} of \url{qbscreen/competence_sensitivity.py}
(seed 20261001; inputs and outputs in \url{calibration_results/phi_bound_test.json}) evaluates
$\Phi_Pk_S/4k_P$ for \nBoundTest\ random spin systems: 0--3 nuclei of spin
$\tfrac12$ or 1 with couplings of 1--300~MHz, exchange $10^{-1}$--$10^4$~MHz of
either sign, axial dipolar tensors of $-10^{-1}$ to $-10^{3}$~MHz with random
axes, random fields up to 0.1~T and rates, and, for about 30\,\% of the small
systems, single-electron dephasing with $T_2^e$ from 1~ns to 10~$\mu$s. The
largest value is \maxBoundRatio.

\section{The search for $F$}
For each $(k_S,k_P)$ on a grid of 15 by 13 rates, the exchange coupling is
scanned on 99 points with $\sinh$ spacing (40~MHz scale), symmetric about zero
and extending to $\max(4000~{\rm MHz},5(k_S+k_P)/2\pi)$, for dipolar scales
$s = 0$, 0.3, 1 and 3 and 25 field directions (the pole and eight azimuths over
$2\pi$ at polar angles of 30, 60 and 90$^\circ$; the yield is invariant under
$B\to-B$). Each of the three largest local maxima in $J$ is refined twice on
21-point subgrids, and the four best distinct candidates are polished by a
Nelder--Mead search in $(J,\theta,\phi,s)$ with $s$ clipped to $[0,3]$ (160
evaluations each). A cell whose maximum lies within 5\,\% of the range limit is
searched again with a range of sixteen times the position of the maximum.
Because the response has narrow resonances in the dipolar scale (at
$k_S = 10^4$, $k_P = 10^{-2}~\mu$s$^{-1}$ and $J\approx0$ it is 0.40, 0.14, 0.68
and 0.04\,\% at $s = 0$, 0.003, 0.01 and 0.03 at the first-pass optimum;
\url{calibration_results/s_resonance_example.json}), every cell is then searched a
second time: $s$ on 41 values (0 and 40 logarithmically spaced from $10^{-3}$ to
3), $J$ on 21 points of the same $\sinh$ grid plus the first-pass optimum, all
25 directions, followed by a polish of the four best candidates
(\url{worst_case_map.supplement}). The cell keeps the larger value.
The second pass raised \nSuppRaised\ cells, by up to a factor of
\suppMaxRatio, without changing any row maximum by more than
\rowMaxChange\,\%. Repeating it with 121 values of $s$ at the cells that
set the thresholds changed $F$ by a factor of at most \convMaxRatio\
(\url{calibration_results/s_convergence.json}). Table~\ref{tab:map}
lists $F$ in percent.

The value $G$ of the main text is computed from these grid values, with the
required rates rounded down to grid rates, and is not a bound on the continuous
rate region. At $\Delta=\offDelta$~eV and $k_{\rm cat}=1$~s$^{-1}$, $k_{\min}$
equals the grid rate $k_P=\offKP~\mu$s$^{-1}$. Nine permitted off-grid values
$k_S=\offKSlo$--$\offKShi~\mu$s$^{-1}$, with $J = 0$ and $\pm0.5$~MHz, $s=0$ and
all 25 field directions, give at most \offMax\,\%, against $G=\offG$\,\%
(\url{competence_sensitivity.offgrid_check};
\url{calibration_results/g_offgrid_check.json}). This is a restricted search at
off-grid points, not a full one.

\begin{table}[h]
\caption{$F(k_S,k_P)$ (\%). Rates in $\mu$s$^{-1}$.}
\label{tab:map}
\centering
\resizebox{\textwidth}{!}{\scriptsize\setlength{\tabcolsep}{2pt}\begin{tabular}{rrrrrrrrrrrrrrrr}
\toprule
$\log k_P \backslash \log k_S$ & -2.0 & -1.5 & -1.0 & -0.5 & 0.0 & 0.5 & 1.0 & 1.5 & 2.0 & 2.5 & 3.0 & 3.5 & 4.0 & 4.5 & 5.0 \\
\midrule
-2.0 & 6.3 & 12 & 14 & 13 & 10 & 7.5 & 6 & 5.2 & 4.6 & 3.7 & 2.3 & 1.3 & 1.3 & 1.3 & 1.3 \\
-1.5 & 2.5 & 6.2 & 11 & 14 & 13 & 9.6 & 7.1 & 5.3 & 3.8 & 2.3 & 1.3 & 1.3 & 1.3 & 1.2 & 1.2 \\
-1.0 & 0.85 & 2.4 & 6 & 11 & 13 & 12 & 8.4 & 5.4 & 2.7 & 1.1 & 1 & 1.1 & 1.1 & 1 & 1 \\
-0.5 & 0.26 & 0.8 & 2.3 & 5.5 & 9.7 & 11 & 9 & 5.2 & 2 & 0.83 & 0.9 & 0.78 & 0.63 & 0.54 & 0.42 \\
0.0 & 0.069 & 0.22 & 0.66 & 1.9 & 4.4 & 7.2 & 7.4 & 4.4 & 1.6 & 0.59 & 0.62 & 0.63 & 0.59 & 0.43 & 0.31 \\
0.5 & 0.014 & 0.044 & 0.14 & 0.41 & 1.1 & 2.5 & 3.6 & 2.7 & 1.1 & 0.37 & 0.34 & 0.37 & 0.35 & 0.26 & 0.15 \\
1.0 & 0.0014 & 0.0044 & 0.014 & 0.043 & 0.13 & 0.34 & 0.66 & 0.75 & 0.48 & 0.2 & 0.07 & 0.084 & 0.088 & 0.077 & 0.045 \\
1.5 & 5.8\text{e-}5 & 1.8\text{e-}4 & 5.8\text{e-}4 & 0.0018 & 0.0056 & 0.017 & 0.043 & 0.083 & 0.096 & 0.062 & 0.032 & 0.049 & 0.055 & 0.042 & 0.02 \\
2.0 & 1.7\text{e-}6 & 5.4\text{e-}6 & 1.7\text{e-}5 & 5.4\text{e-}5 & 1.7\text{e-}4 & 5.3\text{e-}4 & 0.0015 & 0.004 & 0.0077 & 0.014 & 0.018 & 0.02 & 0.023 & 0.015 & 0.0043 \\
2.5 & 1.2\text{e-}7 & 3.8\text{e-}7 & 1.2\text{e-}6 & 3.8\text{e-}6 & 1.2\text{e-}5 & 3.7\text{e-}5 & 1.2\text{e-}4 & 3.5\text{e-}4 & 9.6\text{e-}4 & 0.0021 & 0.0033 & 0.0039 & 0.0036 & 0.0014 & 2.6\text{e-}4 \\
3.0 & 4.8\text{e-}9 & 1.5\text{e-}8 & 4.8\text{e-}8 & 1.5\text{e-}7 & 4.8\text{e-}7 & 1.5\text{e-}6 & 4.8\text{e-}6 & 1.5\text{e-}5 & 4.4\text{e-}5 & 1.1\text{e-}4 & 2.3\text{e-}4 & 2.6\text{e-}4 & 1.4\text{e-}4 & 3.3\text{e-}5 & 4.5\text{e-}6 \\
3.5 & 1.0\text{e-}10 & 3.2\text{e-}10 & 1.0\text{e-}9 & 3.2\text{e-}9 & 1.0\text{e-}8 & 3.2\text{e-}8 & 1.0\text{e-}7 & 3.1\text{e-}7 & 9.5\text{e-}7 & 2.6\text{e-}6 & 5.3\text{e-}6 & 5.4\text{e-}6 & 2.2\text{e-}6 & 3.0\text{e-}7 & 5.7\text{e-}8 \\
4.0 & 1.1\text{e-}12 & 2.6\text{e-}12 & 7.6\text{e-}12 & 2.3\text{e-}11 & 7.4\text{e-}11 & 2.3\text{e-}10 & 7.4\text{e-}10 & 2.3\text{e-}9 & 7.2\text{e-}9 & 2.2\text{e-}8 & 5.8\text{e-}8 & 1.2\text{e-}7 & 1.1\text{e-}7 & 3.3\text{e-}8 & 3.2\text{e-}9 \\
\bottomrule
\end{tabular}
}
\end{table}

\section{Dependence on the spin model}
In the anisotropic model ``aniso'' described below, the full search of
section~S3 (first pass, edge rule and fine-$s$ pass) was repeated on every cell
with $k_P = 10$--$10^{2.5}~\mu$s$^{-1}$ and $k_S\ge k_P$
(\url{model_sensitivity.run_aniso_full};
\url{calibration_results/worst_case_v6_aniso.json}); Table~\ref{tab:aniso}
lists the values, \nAnisoEdge\ of which lie at the edge of the $J$ range. The
thresholds in the main text replace each of these cells by the larger of its
reference and anisotropic values.

The other models were examined only locally: for each onward rate that sets the
thresholds of $G$, the cell that gives the row maximum of $F$ in the reference
model was searched again (\url{qbscreen/model_sensitivity.py}), starting from
the reference optimum and a 39-point $J$ grid at its field direction and dipolar scale, and
polishes the two best candidates. In ``aniso'' N5 and N10 carry the
literature isotropic couplings plus the density-functional spin-dipolar tensors
(Table~\ref{tab:tens}), rotated into the frame of the contact complex in
which the dipolar tensor was computed (0.09~\AA\ root-mean-square deviation
between the relaxed anion and the complex flavin). The complex was built with
methylamine, so the benzylaminium tensors have no orientation in it: the
aminium $\alpha$-protons are kept isotropic (their anisotropy is below 5\,\% of
the isotropic part), and in ``+$^{14}$N'' the aminium nitrogen (spin 1), whose
tensor is nearly axial about the singly occupied orbital, is added with its
axis along the lone pair of the placed amine. ``+2H'' adds the two largest
flavin methyl protons. The relaxation test uses a
reduced pair (N5 isotropic, spin 1; aminium proton 256.9~MHz) searched with the
full procedure, without relaxation and with single-electron dephasing
(Table~\ref{tab:sens}); the reduced-pair relaxation test is superseded by the
direct four-nucleus calculation of section~S5 and is shown for comparison.
Because only the reference row-maximum cell was searched, and only near its
optimum, another cell of the same row may give a larger value; for the
anisotropic model the full search above supersedes these entries.

\begin{table}[h]
\caption{$F(k_S,k_P)$ (\%) in the anisotropic model from the full search.
Rates in $\mu$s$^{-1}$.}
\label{tab:aniso}
\centering
\resizebox{\textwidth}{!}{\scriptsize\setlength{\tabcolsep}{2pt}\begin{tabular}{rrrrrrrrrr}
\toprule
$\log k_P \backslash \log k_S$ & 1.0 & 1.5 & 2.0 & 2.5 & 3.0 & 3.5 & 4.0 & 4.5 & 5.0 \\
\midrule
1.0 & 1.7 & 2.3 & 1.8 & 0.99 & 0.52 & 0.28 & 0.18 & 0.14 & 0.067 \\
1.5 &  & 0.23 & 0.28 & 0.21 & 0.13 & 0.068 & 0.026 & 0.017 & 0.0082 \\
2.0 &  &  & 0.022 & 0.027 & 0.021 & 0.011 & 0.0049 & 0.0019 & 5.1\text{e-}4 \\
2.5 &  &  &  & 0.0021 & 0.003 & 0.0025 & 0.0014 & 3.7\text{e-}4 & 6.4\text{e-}5 \\
\bottomrule
\end{tabular}
}
\end{table}

\begin{table}[h]
\caption{Searched $F$ (\%) at the threshold-setting cells in the reference and
richer models, and in the reduced pair with and without electron dephasing.}
\label{tab:sens}
\centering
\scriptsize\setlength{\tabcolsep}{3pt}\begin{tabular}{rrrrrrrrrrr}
\toprule
$\log k_P$ & $\log k_S$ & reference & aniso & +$^{14}$N & +2H & red.\ coh. & 10 ns & 100 ns & 1 $\mu$s & 10 $\mu$s \\
\midrule
1.0 & 1.5 & 0.75 & 2.3 & 1 & 1.6 & 0.51 & 0.21 & 0.22 & 0.41 & 0.5 \\
1.5 & 2.0 & 0.096 & 0.28 & 0.16 & 0.2 & 0.086 & 0.052 & 0.038 & 0.076 & 0.085 \\
2.0 & 4.0 & 0.023 & 0.0035 & 0.0033 & 0.0036 & 0.0033 & 0.016 & 0.0028 & 0.0032 & 0.0032 \\
2.5 & 3.5 & 0.0039 & 0.0023 & 0.0015 & 0.0022 & 7.8\text{e-}4 & 0.0012 & 7.6\text{e-}4 & 7.8\text{e-}4 & 7.8\text{e-}4 \\
\bottomrule
\end{tabular}

\end{table}

\section{Electron spin relaxation in the four-nucleus model}
The four-nucleus spin space has dimension 144, too large for the dense
Liouville-space solution. With $\gamma = 2/T_2^e$, single-electron dephasing
$\gamma\sum_e(S_{ez}\rho S_{ez}-\rho/4)$ is split into a decay $-(\gamma/2)\rho$,
absorbed into $K = 2\pi H - \tfrac i2[k_SP_S + (k_P+\gamma/2)I]$, and a recycling
term $\gamma\sum_eS_{ez}\rho S_{ez}$. The integral $X=\int_0^\infty\rho\,dt$ then
solves $KX - XK^\dagger = -i\rho_0 - i\gamma\sum_eS_{ez}XS_{ez}$, which is iterated
by GMRES\cite{Saad1986} with the left side inverted by a triangular Sylvester solve in the Schur
basis of $K$ (\url{product_yield.yields_t2e}). On twelve random systems small
enough for the dense solution, including anisotropic hyperfine tensors, the two
agree to a relative $1.3\times10^{-13}$. For every cell with $k_S\ge k_P$ on the
threshold rows, the first-pass search of section~S3 was repeated with
$T_2^e = 10$~ns, 100~ns and 1~$\mu$s, with the same edge rule for the $J$ range
(Table~\ref{tab:relax}; \url{calibration_results/worst_case_v6_t2e.json}).
Nineteen searches (14 at $T_2^e = 10$~ns, 5 at 100~ns) first peaked at the edge
of the $J$ range and were repeated with the range widened sixteenfold; one of them
($k_S = 10^{4.5}$, $k_P = 10~\mu$s$^{-1}$, $T_2^e = 100$~ns) still peaks at the
edge because the response approaches a plateau as $|J|\to\infty$: it is
constant to five figures for $|J| = 10^6$--$10^8$~MHz and equal to the stored
value (\url{calibration_results/t2e_plateau_check.json}).

\begin{table}[h]
\caption{Searched $F$ (\%) without and with single-electron dephasing in the
four-nucleus model. Rates in $\mu$s$^{-1}$.}
\label{tab:relax}
\centering
\small\begin{tabular}{rrrrrr}
\toprule
$\log k_P$ & $\log k_S$ & coherent & 10 ns & 100 ns & 1000 ns \\
\midrule
1.0 & 1.0 & 0.66 & 0.075 & 0.22 & 0.54 \\
1.0 & 1.5 & 0.75 & 0.12 & 0.29 & 0.63 \\
1.0 & 2.0 & 0.48 & 0.18 & 0.26 & 0.41 \\
1.0 & 2.5 & 0.2 & 0.21 & 0.18 & 0.18 \\
1.0 & 3.0 & 0.07 & 0.22 & 0.21 & 0.067 \\
1.0 & 3.5 & 0.084 & 0.22 & 0.21 & 0.073 \\
1.0 & 4.0 & 0.088 & 0.22 & 0.21 & 0.076 \\
1.0 & 4.5 & 0.077 & 0.22 & 0.21 & 0.065 \\
1.0 & 5.0 & 0.045 & 0.22 & 0.21 & 0.04 \\
1.5 & 1.5 & 0.083 & 0.024 & 0.046 & 0.077 \\
1.5 & 2.0 & 0.096 & 0.042 & 0.059 & 0.09 \\
1.5 & 2.5 & 0.062 & 0.063 & 0.045 & 0.059 \\
1.5 & 3.0 & 0.032 & 0.072 & 0.028 & 0.03 \\
1.5 & 3.5 & 0.049 & 0.075 & 0.035 & 0.047 \\
1.5 & 4.0 & 0.055 & 0.076 & 0.04 & 0.053 \\
1.5 & 4.5 & 0.042 & 0.077 & 0.033 & 0.041 \\
1.5 & 5.0 & 0.02 & 0.077 & 0.015 & 0.019 \\
2.0 & 2.0 & 0.0077 & 0.0067 & 0.0068 & 0.0076 \\
2.0 & 2.5 & 0.014 & 0.011 & 0.01 & 0.014 \\
2.0 & 3.0 & 0.018 & 0.013 & 0.014 & 0.017 \\
2.0 & 3.5 & 0.02 & 0.015 & 0.016 & 0.019 \\
2.0 & 4.0 & 0.023 & 0.015 & 0.018 & 0.022 \\
2.0 & 4.5 & 0.015 & 0.015 & 0.012 & 0.014 \\
2.0 & 5.0 & 0.0043 & 0.015 & 0.0033 & 0.0042 \\
2.5 & 2.5 & 0.0021 & 0.0016 & 0.002 & 0.0021 \\
2.5 & 3.0 & 0.0033 & 0.0024 & 0.0031 & 0.0032 \\
2.5 & 3.5 & 0.0039 & 0.0023 & 0.0036 & 0.0039 \\
2.5 & 4.0 & 0.0036 & 0.0017 & 0.0031 & 0.0035 \\
2.5 & 4.5 & 0.0014 & 0.0013 & 0.0012 & 0.0013 \\
2.5 & 5.0 & 2.6\text{e-}4 & 0.0012 & 2.2\text{e-}4 & 2.5\text{e-}4 \\
\bottomrule
\end{tabular}

\end{table}

\section{Hyperfine couplings}
Fermi-contact couplings (B3LYP/def2-TZVP at GFN2-xTB geometries), all nuclei
with $|a_{\rm iso}|\ge1$~MHz (Table~\ref{tab:hfc}), and the spin-dipolar tensors
of the nitrogens and the largest protons (Table~\ref{tab:tens}), computed from
the spin density with the operator-derivative integrals
$\langle\mu|\partial_a\partial_b|\mathbf r-\mathbf R|^{-1}|\nu\rangle$; their
trace reproduces $-4\pi\rho_s(\mathbf R)$. The computed flavin N5 and N10
isotropic couplings are within 16\,\% and 13\,\% of the literature values used in
the main text,\cite{Lee2014} and the N5 tensor is nearly axial with its large
component normal to the ring. For 3-methylindole$^{\bullet+}$ the two large
methyl couplings stand in for the tryptophan $\beta$-protons; their values
depend on the side-chain conformation.

\begin{table}[h]
\caption{Isotropic hyperfine couplings.}
\label{tab:hfc}
\centering
\small\begin{tabular}{llrr}
\toprule
radical & nucleus (atom index) & $a_{\rm iso}$ (MHz) & $a_{\rm iso}$ ($\mu$T) \\
\midrule
lumiflavin$^{\bullet-}$ & H29 & +17.9 & +639 \\
 & H30 & +17.8 & +636 \\
 & N4 & +12.4 & +441 \\
 & H25 & +11.3 & +403 \\
 & H26 & +11.3 & +402 \\
 & H22 & -10.5 & -374 \\
 & H20 & -5.3 & -190 \\
 & H21 & -5.3 & -190 \\
 & N13 & +4.6 & +165 \\
 & H27 & +1.4 & +50 \\
 & N8 & -1.0 & -36 \\
methylaminium$^{\bullet+}$ & H2 & +363.1 & +12955 \\
 & H6 & +82.7 & +2951 \\
 & H5 & +81.9 & +2923 \\
 & H3 & -59.3 & -2116 \\
 & H4 & -59.3 & -2115 \\
 & N1 & +30.5 & +1090 \\
benzylaminium$^{\bullet+}$ & H10 & +256.9 & +9167 \\
 & H11 & +196.2 & +7001 \\
 & N0 & +24.8 & +885 \\
 & H8 & -23.0 & -821 \\
 & H9 & -20.7 & -737 \\
 & H14 & -12.3 & -439 \\
 & H16 & -7.8 & -278 \\
 & H13 & -3.2 & -115 \\
 & H15 & +2.4 & +85 \\
3-methylindole$^{\bullet+}$ & H12 & +60.1 & +2143 \\
 & H11 & +55.5 & +1980 \\
 & H18 & -14.3 & -511 \\
 & H13 & -13.1 & -467 \\
 & H14 & -13.1 & -466 \\
 & H16 & -9.8 & -351 \\
 & H15 & -4.9 & -175 \\
 & N3 & +4.8 & +172 \\
 & H17 & +2.7 & +95 \\
\bottomrule
\end{tabular}

\end{table}

\begin{table}[h]
\caption{Spin-dipolar hyperfine tensors.}
\label{tab:tens}
\centering
\small\begin{tabular}{llrl}
\toprule
radical & nucleus & $a_{\rm iso}$ (MHz) & principal values of $\mathsf A_{\rm dip}$ (MHz) \\
\midrule
lumiflavin$^{\bullet-}$ & H29 & +17.9 & -1.2, -1.0, +2.2 \\
 & H30 & +17.8 & -1.2, -1.0, +2.2 \\
 & N4 & +12.4 & -17.0, -16.7, +33.7 \\
 & H25 & +11.3 & -1.5, -0.9, +2.4 \\
 & H26 & +11.3 & -1.5, -0.9, +2.4 \\
 & N13 & +4.6 & -6.0, -5.7, +11.7 \\
 & N8 & -1.0 & -0.2, +0.1, +0.1 \\
 & N11 & -0.2 & -0.6, -0.4, +1.1 \\
benzylaminium$^{\bullet+}$ & H10 & +256.9 & -8.0, -4.2, +12.2 \\
 & H11 & +196.2 & -7.9, -4.2, +12.1 \\
 & N0 & +24.8 & -18.8, -18.6, +37.4 \\
 & H8 & -23.0 & -21.8, -6.0, +27.8 \\
 & H9 & -20.7 & -21.4, -5.4, +26.8 \\
\bottomrule
\end{tabular}

\end{table}

\section{Dipolar tensor}
The tensor is the point-dipole tensor summed over pairs of atoms weighted by
their Löwdin\cite{Lowdin1950} spin populations (UKS B3LYP/def2-SVP on each radical at its position
in the complex), for methylamine on the flavin si face with N5--N = 3.5~\AA\
(Table~\ref{tab:dip}).

\begin{table}[h]
\caption{Dipolar coupling of the contact pair.}
\label{tab:dip}
\centering
\begin{tabular}{lrrl}
\toprule
model & $D$ (MHz) & $E$ (MHz) & principal values of $\mathsf T$ (MHz) \\
\midrule
spin-population tensor & -843 & -26 & -1124, 511, 613 \\
point dipole at spin centroids & -1642 & 0 & -- \\
point dipole at N5 and N & -1817 & 0 & -- \\
\bottomrule
\end{tabular}

\end{table}

\section{Electronic coupling}
Fragment-orbital couplings\cite{Valeev2006} $t = \langle{\rm LUMO_{Fl}}|F|{\rm HOMO_{am}}\rangle$
(B3LYP/def2-SVP) for methylamine and benzylamine stacked on the flavin si face
(five orientations) or approaching its edge (three orientations), at N5--N
distances of 3.0--6.0~\AA. Geometries with any intermolecular contact below
2.0~\AA\ are excluded (Table~\ref{tab:t}). The edge approach of benzylamine
clashes at every distance below 6~\AA\ and is not used.

\begin{table}[h]
\caption{Coupling scan. ``valid'': non-clashing geometries of those computed.}
\label{tab:t}
\centering
\small\begin{tabular}{llrrrrr}
\toprule
amine & approach & $r_{\rm N5N}$ (\AA) & valid & median $|t|$ (meV) & max $|t|$ (meV) & min contact (\AA) \\
\midrule
methylamine & face & 3.0 & 5/5 & 184.6 & 311.6 & 2.32 \\
methylamine & face & 3.5 & 5/5 & 52.8 & 105.6 & 2.72 \\
methylamine & face & 4.0 & 5/5 & 24.6 & 37.9 & 3.12 \\
methylamine & face & 4.5 & 5/5 & 5.1 & 31.1 & 3.53 \\
methylamine & face & 5.0 & 5/5 & 1.9 & 14.7 & 3.96 \\
methylamine & face & 5.5 & 5/5 & 1.0 & 4.2 & 4.40 \\
methylamine & face & 6.0 & 5/5 & 0.3 & 1.1 & 4.85 \\
methylamine & edge & 3.0 & 0/3 & -- & -- & 0.77 \\
methylamine & edge & 3.5 & 0/3 & -- & -- & 0.93 \\
methylamine & edge & 4.0 & 0/3 & -- & -- & 1.29 \\
methylamine & edge & 4.5 & 2/3 & 97.0 & 170.8 & 1.71 \\
methylamine & edge & 5.0 & 3/3 & 23.1 & 81.0 & 2.17 \\
methylamine & edge & 5.5 & 3/3 & 8.6 & 35.0 & 2.65 \\
methylamine & edge & 6.0 & 3/3 & 3.2 & 12.7 & 3.13 \\
benzylamine & face & 3.0 & 0/5 & -- & -- & 0.61 \\
benzylamine & face & 3.5 & 2/5 & 47.1 & 52.3 & 0.89 \\
benzylamine & face & 4.0 & 2/5 & 5.2 & 7.9 & 1.31 \\
benzylamine & face & 4.5 & 3/5 & 6.1 & 35.2 & 1.77 \\
benzylamine & face & 5.0 & 5/5 & 4.2 & 21.1 & 2.25 \\
benzylamine & face & 5.5 & 5/5 & 1.0 & 8.3 & 2.73 \\
benzylamine & face & 6.0 & 5/5 & 0.9 & 2.0 & 3.22 \\
benzylamine & edge & 3.0 & 0/3 & -- & -- & 1.00 \\
benzylamine & edge & 3.5 & 0/3 & -- & -- & 1.00 \\
benzylamine & edge & 4.0 & 0/3 & -- & -- & 0.92 \\
benzylamine & edge & 4.5 & 0/3 & -- & -- & 0.82 \\
benzylamine & edge & 5.0 & 0/3 & -- & -- & 1.00 \\
benzylamine & edge & 5.5 & 0/3 & -- & -- & 1.36 \\
benzylamine & edge & 6.0 & 1/3 & 0.4 & 0.4 & 1.78 \\
\bottomrule
\end{tabular}

\end{table}

\section{Density-functional driving force}
The main text uses measured potentials. For comparison, Table~\ref{tab:pot}
gives adiabatic one-electron potentials in water ($\varepsilon = 78.4$, absolute
SHE potential 4.44~V\cite{Jonsson1996}) from the same density-functional
protocol, against the measured values for a dialkylamine\cite{Jonsson1996} and
for free FMN at pH~7.\cite{Mayhew1999} The flavin comparison is between
different species (lumiflavin anion radical versus the pH-7 FMN
oxidized/semiquinone couple) and indicates only the order of the error.
Table~\ref{tab:ddft} gives the raw driving force in low-dielectric continua with
a contact Coulomb term $-e^2/(4\pi\varepsilon_0\varepsilon r)$. No zero-point or
thermal corrections are included, and none of these values enters the main
text.

\begin{table}[h]
\caption{One-electron potentials in water.}
\label{tab:pot}
\centering
\begin{tabular}{lrrr}
\toprule
couple & DFT (V vs NHE) & experiment & error (V) \\
\midrule
Me$_2$NH$^{\bullet+}$/Me$_2$NH & +0.73 & +1.300 & -0.57 \\
MeNH$_2^{\bullet+}$/MeNH$_2$ & +1.26 & -- & -- \\
BnNH$_2^{\bullet+}$/BnNH$_2$ & +1.45 & -- & -- \\
Fl/Fl$^{\bullet-}$ (exp.: FMN, pH 7) & -1.16 & -0.313 & -0.84 \\
\bottomrule
\end{tabular}

\end{table}

\begin{table}[h]
\caption{Raw DFT driving force of the contact pair.}
\label{tab:ddft}
\centering
\begin{tabular}{llrrrr}
\toprule
continuum & amine & separated (eV) & $r=4$ \AA & 5 \AA & 6 \AA \\
\midrule
$\varepsilon=4$ & methylamine & 3.55 & 2.65 & 2.83 & 2.95 \\
$\varepsilon=4$ & benzylamine & 3.59 & 2.69 & 2.87 & 2.99 \\
$\varepsilon=10$ & methylamine & 2.83 & 2.47 & 2.55 & 2.59 \\
$\varepsilon=10$ & benzylamine & 2.97 & 2.61 & 2.69 & 2.73 \\
\bottomrule
\end{tabular}

\end{table}

\section{The driving force in the active site}
The main text places the pair in water at
$\Delta_w = E^\circ_{\rm amine} - E_{\rm fl} - e^2/(4\pi\varepsilon_0\cdot78.4\,r)$,
$r = 3.5$--6~\AA. Moving the pair into a medium of effective dielectric constant
$\varepsilon$ changes the Coulomb term by
$-e^2/(4\pi\varepsilon_0r)\,(1/\varepsilon-1/78.4)$ and the solvation of the ions,
which is not known. Two continuum limits are evaluated for $\varepsilon = 4$ and 10.

\emph{One-sided.} Only the aminium radical cation is desolvated, with a Born
term $e^2/(8\pi\varepsilon_0a)\,(1/\varepsilon-1/78.4)$; the flavin anion keeps its
aqueous solvation. The radius $a = \BornRadius$~\AA\ reproduces, by the Born
formula in water, the solvation-energy difference of 244~kJ\,mol$^{-1}$ between
dimethylamine and its radical cation.\cite{Jonsson1996} The change of $\Delta$
is $\MedOneLo$ to $+\MedOneHi$~eV over $\varepsilon$ and $r$; it is negative
because the stronger attraction of the pair outweighs the desolvation of one ion,
and it is the most favourable case for a small $\Delta$ that the continuum
picture allows.

\emph{Two-ion.} Both ions are moved with the density-functional continuum
protocol of the previous section, using the separated-ion energies at
$\varepsilon = 4$ and 10 (Table~\ref{tab:ddft}) against those at
$\varepsilon = 78.4$ (Table~\ref{tab:pot}), for methylamine and benzylamine with
lumiflavin. The change of $\Delta$, including the Coulomb term, is $+\MedTwoLo$ to
$+\MedTwoHi$~eV.

\emph{Substrate-induced shift.} Without substrate the oxidized/semiquinone and
semiquinone/reduced midpoints of the MAO~A flavin are $\EflMAOA$ and
$\EsqredA$~mV,\cite{Sablin2001} so the two-electron midpoint is $\EmFreeA$~mV.
With substrate the full-reduction midpoint shifts by up to $+\SubShift$~V and no
semiquinone is observed.\cite{Sablin2001} If the absence of semiquinone reflects
equilibrium, the first one-electron potential lies at or below the two-electron
midpoint, at most $\EmFreeA+500$~mV${}=+\EoneCap$~V, a shift of the one-electron
couple of at most $\EoneShiftCap$~V. The main text uses the full
$+\SubShift$~V, which lowers $\Delta$ more than this reading allows. All three
estimates are computed in \url{qbscreen/make_numbers_v6.py}.

\section{Cryptochrome}
Orientation-averaged, perpendicular and parallel response of the cryptochrome
pair over the exchange prefactor $J_0$ (both signs of $J_{\rm EH}$;
$J = -2J_{\rm EH}$, $J_{\rm EH} = J_0e^{-\beta r}$, $\beta = 14$~nm$^{-1}$,
$r = 1.90$~nm\cite{Efimova2008}), $k_S = k_P = 1~\mu$s$^{-1}$, coherent
(Table~\ref{tab:cry}).

\begin{table}[h]
\caption{Cryptochrome response (\%).}
\label{tab:cry}
\centering
\small\begin{tabular}{rrrrr}
\toprule
$J_0$ ($\mu$T) & $J$ (MHz) & powder (\%) & $\perp$ (\%) & $\parallel$ (\%) \\
\midrule
8.00\text{e}12 & -1.26 & +0.417 & +0.494 & +0.027 \\
8.00\text{e}12 & +1.26 & +0.386 & +0.464 & -0.023 \\
1.01\text{e}13 & -1.58 & +0.395 & +0.469 & +0.015 \\
1.01\text{e}13 & +1.58 & +0.384 & +0.457 & -0.030 \\
1.27\text{e}13 & -1.99 & +0.371 & +0.454 & -0.005 \\
1.27\text{e}13 & +1.99 & +0.392 & +0.461 & -0.034 \\
1.60\text{e}13 & -2.51 & +0.403 & +0.498 & +0.005 \\
1.60\text{e}13 & +2.51 & +0.413 & +0.491 & -0.029 \\
2.01\text{e}13 & -3.16 & +0.470 & +0.561 & +0.021 \\
2.01\text{e}13 & +3.16 & +0.435 & +0.536 & -0.020 \\
2.53\text{e}13 & -3.98 & +0.573 & +0.681 & +0.027 \\
2.53\text{e}13 & +3.98 & +0.444 & +0.580 & -0.012 \\
3.18\text{e}13 & -5.01 & +0.634 & +0.751 & +0.042 \\
3.18\text{e}13 & +5.01 & +0.404 & +0.483 & +0.007 \\
4.01\text{e}13 & -6.30 & +0.559 & +0.622 & +0.034 \\
4.01\text{e}13 & +6.30 & +0.334 & +0.407 & +0.025 \\
5.05\text{e}13 & -7.93 & +0.383 & +0.420 & +0.061 \\
5.05\text{e}13 & +7.93 & +0.242 & +0.320 & -0.005 \\
6.35\text{e}13 & -9.99 & +0.441 & +0.534 & +0.033 \\
6.35\text{e}13 & +9.99 & +0.242 & +0.299 & +0.049 \\
8.00\text{e}13 & -12.57 & +0.492 & +0.589 & +0.015 \\
8.00\text{e}13 & +12.57 & +0.168 & +0.234 & +0.021 \\
1.01\text{e}14 & -15.83 & +0.282 & +0.372 & +0.024 \\
1.01\text{e}14 & +15.83 & +0.153 & +0.174 & +0.081 \\
1.27\text{e}14 & -19.93 & +0.085 & +0.115 & -0.060 \\
1.27\text{e}14 & +19.93 & +0.150 & +0.131 & +0.126 \\
1.60\text{e}14 & -25.09 & +0.061 & +0.078 & +0.020 \\
1.60\text{e}14 & +25.09 & +0.185 & +0.119 & +0.214 \\
2.01\text{e}14 & -31.58 & +0.044 & +0.049 & +0.023 \\
2.01\text{e}14 & +31.58 & +0.055 & +0.078 & +0.063 \\
2.53\text{e}14 & -39.76 & -0.009 & -0.004 & -0.013 \\
2.53\text{e}14 & +39.76 & +0.036 & +0.041 & +0.019 \\
3.18\text{e}14 & -50.05 & +0.015 & +0.017 & +0.009 \\
3.18\text{e}14 & +50.05 & +0.068 & +0.063 & +0.041 \\
4.01\text{e}14 & -63.01 & +0.010 & +0.011 & +0.007 \\
4.01\text{e}14 & +63.01 & +0.052 & +0.045 & +0.031 \\
5.05\text{e}14 & -79.33 & +0.004 & +0.005 & +0.003 \\
5.05\text{e}14 & +79.33 & +0.039 & +0.049 & +0.015 \\
6.35\text{e}14 & -99.87 & +0.003 & +0.003 & +0.002 \\
6.35\text{e}14 & +99.87 & +0.034 & +0.036 & +0.019 \\
8.00\text{e}14 & -125.73 & +0.001 & +0.002 & +0.001 \\
8.00\text{e}14 & +125.73 & +0.011 & +0.017 & +0.002 \\
\bottomrule
\end{tabular}

\end{table}

\bibliographystyle{unsrt}
\bibliography{refs}

\end{document}
